\section{The Imported Action and Its Application}
\label{sec:entropic_action}

For this application we use a proposed specialization of the PM operational action \cite{kate_meyer_persistence_2026}. It models the cost of maintaining required distinctions by a density
\begin{equation}
 \mathcal D=\xi\,\mathbf K\,f,\qquad S_\Phi=\int\mathcal D\,d\mu,
\end{equation}
where $\xi$ is a specified primitive cost, $\mathbf K$ the maintained workload or complexity density, and $f$ its operation rate. The state space, counted events, units and measure must be supplied for an application. Information capacity, description length and executed work are not interchangeable without that map. The multiplicative density is a modeling contract, not a proved consequence of statistical independence from coding or thermodynamic theorems.

A resource budget bounds the product and hence bounds $f$ when $\xi$ and $\mathbf K$ are positive (Lean: \nolinkurl{PMAction.complexity_speed_bound}). It does not fix the factors individually. The economical selection principle is part of the application contract; this paper does not rederive a universal minimum-entropy-production theorem. A closed vacuum also cannot be treated as shedding entropy to an external bath without specifying the internal statistical description.

\subsection{Least Action as Survivorship}
\label{sec:least_action}
For a finite history family with costs $S_i$, consider the supplied weights
\begin{equation}
 P_\beta(i)=\frac{e^{-\beta S_i}}{\sum_j e^{-\beta S_j}}.
 \label{eq:survival_prob}
\end{equation}
Lean proves positivity and normalization, multiplicative weights for additive costs, and the exact ratio
\begin{equation}
 \frac{P_\beta(j)}{P_\beta(i)}=e^{-\beta(S_j-S_i)}.
\end{equation}
For finite real costs and tolerance, every weight is strictly positive. Thus suppression is not exact deletion. For a fixed finite family, \nolinkurl{PMFramework.Gibbs.probability_limit} proves that the normalized distribution tends to the uniform distribution on its minimum-cost histories as $\beta\to+\infty$. A thermodynamic or long-observation application must justify its limit and effective inverse tolerance.

Stationarity alone is weaker than a minimum: \nolinkurl{PMAction.stationary_not_minimal} supplies a counterexample. A nonnegative dissipative functional also cannot be equated directly with an unrestricted Lorentzian quadratic form (\nolinkurl{lorentzian_cost_not_nonnegative}). The relation to physical time evolution therefore requires an additional construction, discussed next.


\subsection{A Constructive Carrier--Event Cost Connection}
The counting law can now be applied to explicit substrate carriers. \nolinkurl{CarrierEvents}
uses Paper I's minimal doubled vectors in each nontrivial $D_4$ glue class; every class
has eight states, with no minimal vector omitted by the finite encoding. The reflection
in $e_0-e_1$ permutes each class reversibly and moves four of its eight states.
\nolinkurl{SubstrateReflection} connects that coordinate reflection to an integral $E_8$
root reflection preserving the actual quadratic form and glue labels, including its
restriction through Paper I's $D_4\oplus D_4$ embedding.

With uniform carrier sampling, a changed-state event has probability $1/2$
and cost $c_0/2$. The uniform law is not chosen: under the framework's cost-proportional hazard, it is the
sampling law that survives selection among the candidates (Persistence Mechanics, \S4; framework Lean \nolinkurl{PMFramework.LawSelection.selected_uniform_average}),
and on each carrier class the $D_4$ Weyl reflections act transitively (Paper~I,
\nolinkurl{PMArch.ChannelSymmetry}); \nolinkurl{PMAction.SelectedSampling.selected_primitive_cost}. More generally, a permutation $P$ on $n>0$ carrier labels has
\begin{equation}
 p_{\rm changed}=1-\frac{\operatorname{Re}\operatorname{Tr}P}{n},\qquad
 C_{\rm changed}=c_0p_{\rm changed}
\end{equation}
(\nolinkurl{CarrierCost}). The trace score here follows from counting fixed labels.
This does not uniquely extend a permutation event to a continuous phase readout.
Nor are eight distinct geometric vectors eight independent vectors: Lean verifies
that the carrier vectors in four coordinates are linearly dependent. A complex
amplitude space over their labels is an additional construction, not that cardinality identity.

\paragraph{Phase-sensitive comparison.}
For a specified unitary transport $U$ and a reference branch, take the minus-port
amplitude of a prepared basis label $a$ to be $(e_a-Ue_a)/2$. Under the stated
squared-amplitude readout, Lean derives its probability and the uniform mean
\begin{equation}
 p_{\rm flag}(a)=\frac{1-\operatorname{Re}U_{aa}}2,\qquad
 \overline p_{\rm flag}=\frac12\left(1-\frac{\operatorname{Re}\operatorname{Tr}U}{n}\right).
\end{equation}
The two output probabilities are nonnegative and sum to one
(\nolinkurl{InterferenceReadout}). If a flag triggers $m$ primitive repair events,
weighted finite-event accounting gives $C=c_0m\,\overline p_{\rm flag}$, without an
independence assumption (\nolinkurl{EventAccounting}, \nolinkurl{CarrierResponse}).
For the carrier reflection this flag probability is $1/4$, whereas its direct
changed-state probability is $1/2$. The operational event must therefore be specified.

For a scalar phase $U=e^{i\theta}$ the repair cost is $(c_0m/2)(1-\cos\theta)$.
A phase-independent required comparison cost remains a baseline $B$; subtracting
the zero-phase reference gives that excess cost, rather than declaring comparison
free. Reference preparation, squared-amplitude readout and the
physical repair multiplicity remain explicit inputs; uniform sampling is the selected law. This is not a derivation of
the Born rule or a unique choice of the remaining inputs from persistence.
